\section*{第 \ref{ch:g7:earthquakes} 章——地震}

\begin{solution}{exo:g7:earthquakes:1}
(a) \omterm{def:g7:earthquakes:focus}{震源}；(b) \omterm{def:g7:earthquakes:intensity}{烈度}；(c) \omterm{def:g7:earthquakes:fault}{断层}；(d) \omterm{def:g7:earthquakes:focus}{震中}；(e) \omterm{def:g7:earthquakes:magnitude}{震级}；(f)
\omterm{def:g7:earthquakes:seismometer}{地震仪}。
\end{solution}

\begin{solution}{exo:g7:earthquakes:2}
(a) 就在村庄处：\omterm{def:g7:earthquakes:focus}{震中}是地表上
正对\omterm{def:g7:earthquakes:focus}{震源}的那一点。(b) \omterm{def:g7:earthquakes:focus}{震源}在正下方 \qty{30}{km} 处。(c) 通常
较弱：震动随距离减弱。（这也取决于地面：
在松软的泥土上，第二个村庄仍可能强烈震动。）
\end{solution}

\begin{solution}{exo:g7:earthquakes:3}
(a) \omterm{def:g7:earthquakes:magnitude}{震级}（带一位小数的数）。(b) \omterm{def:g7:earthquakes:intensity}{烈度}（在某一地点所见的
情况）。(c) \omterm{def:g7:earthquakes:intensity}{烈度}（罗马数字，两地不同）。
(d) \omterm{def:g7:earthquakes:magnitude}{震级}：表示释放的能量，\omterm{def:g7:earthquakes:magnitude}{震级}每增加一级。
\end{solution}

\begin{solution}{exo:g7:earthquakes:4}
(a) 错误：科学家无法预测\omterm{def:g4:volcanoes-earthquakes-first-look:earthquake}{地震}，他们只知道哪些
\omterm{def:g7:earthquakes:fault}{断层}将来会再次破裂。(b) 错误：在开阔\omterm{def:g2:water-on-earth:ocean}{海洋}中，\omterm{def:g7:earthquakes:tsunami}{海啸}
高度很少超过一米，靠近\omterm{def:g1:landscapes-around-us:water}{海岸}时才会增高。(c) 正确。(d)
正确：\omterm{def:g7:earthquakes:magnitude}{震级} 6 级或以上的\omterm{def:g4:volcanoes-earthquakes-first-look:earthquake}{地震}中，92~\% 发生在
距\omterm{def:g7:plate-tectonics:plate-boundary}{板块边界}不到 \qty{300}{km} 处。
\end{solution}

\begin{solution}{exo:g7:earthquakes:5}
(a) 每年约 127 次\omterm{def:g7:earthquakes:magnitude}{震级}为 6 至 6.9 级的\omterm{def:g4:volcanoes-earthquakes-first-look:earthquake}{地震}；约 13 次\omterm{def:g7:earthquakes:magnitude}{震级}为 7 至
7.9 级的\omterm{def:g4:volcanoes-earthquakes-first-look:earthquake}{地震}。(b) $126.7 \div 13.12 \approx 9.7$：约少 10 倍。(c)
$13.12 + 0.76 + 0.04 = 13.92$：每年约 \textbf{14} 次。(d)
$0.04 \times 50 = 2$：50 年内两次（苏门答腊 2004 年和日本 2011 年）。
\end{solution}

\begin{solution}{exo:g7:earthquakes:6}
(a) 一级：约 32 倍。(b) 两级：$32 \times 32 = 1024$，约
1000 倍。(c) 每升一级，波就大 10 倍：$10 \times 10 =
100$ 倍。
\end{solution}

\begin{solution}{exo:g7:earthquakes:7}
(a) $8 \times 5 = \qty{40}{km}$。(b) $8 \times 25 = \qty{200}{km}$。(c)
$120 \div 8 = \qty{15}{s}$。(d) 一个台站只能给出距离，不能给出
方向：到该台站距离相同的所有点，构成一个以台站为圆心的圆，
\omterm{def:g7:earthquakes:focus}{震中}可能位于该圆上的任何一点。
\end{solution}

\begin{solution}{exo:g7:earthquakes:8}
(a) V。(b) II。(c) VIII。(d) \omterm{def:g7:earthquakes:intensity}{烈度}因地而异：
报告 (c) 来自靠近\omterm{def:g7:earthquakes:focus}{震中}的地方，或建于松软
地面上；报告 (a) 来自较远处，或位于坚硬\omterm{def:g3:stones-pebbles-sand:rock}{岩石}上。
这次\omterm{def:g4:volcanoes-earthquakes-first-look:earthquake}{地震}只有一个\omterm{def:g7:earthquakes:magnitude}{震级}。
\end{solution}

\begin{solution}{exo:g7:earthquakes:9}
(a) $7031 - 5561 = 1470$。(b) $1470 \div 7031 \approx 0.209$：约
\textbf{21~\%}。(c) 下沉的板块深入\omterm{def:g6:inside-the-earth:layers}{地幔}。它
仍足够坚硬而能断裂，因此\omterm{def:g4:volcanoes-earthquakes-first-look:earthquake}{地震}可在其内部
数百千米深处发生。
\end{solution}

\begin{solution}{exo:g7:earthquakes:10}
(a) $2500 \div 800 = 3.125$ 小时：3 小时加 $0.125 \times 60 = 7.5$
分钟，约 3 小时 8 分钟。(b) 3 小时 8 分钟 $-$ 2 小时 17 分钟 $=$
晚 51 分钟。(c) \omterm{def:g7:earthquakes:tsunami}{海啸}波用了数小时才到达，而海底大\omterm{def:g4:volcanoes-earthquakes-first-look:earthquake}{地震}
通常几分钟内即可定位：发出预警，本可在\omterm{def:g7:earthquakes:tsunami}{海啸}到来很久以前
让人们转移到高处。
\end{solution}

\begin{solution}{exo:g7:earthquakes:11}
(a) $400 \div 800 = 0.5$ 小时：30 分钟。(b) $20 \div 40 = 0.5$ 小时：
30 分钟。(c) 60 分钟。(d) 最后 \qty{20}{km} 所用的时间
与最初 \qty{400}{km} 相同：波在浅水中急剧减慢。(e)
波前减速，而后方的水仍快速涌来：水
堆积起来，波变得更高。
\end{solution}

\begin{solution}{exo:g7:earthquakes:12}
(a) 三级：$32 \times 32 \times 32 = 32\,768$，约
\textbf{33\,000} 倍。(b) $32\,768 \div 127 \approx 258$：约
相当于全球所有 6 级\omterm{def:g4:volcanoes-earthquakes-first-look:earthquake}{地震} 260 年的总量。(c) 几乎全部
能量由少数几次最大的\omterm{def:g4:volcanoes-earthquakes-first-look:earthquake}{地震}释放：一次
9 级\omterm{def:g4:volcanoes-earthquakes-first-look:earthquake}{地震}的能量，超过数百年的 6 级\omterm{def:g4:volcanoes-earthquakes-first-look:earthquake}{地震}。
\end{solution}

\begin{solution}{pb:g7:earthquakes:1}
\textbf{1.} A：$45 - 25 = 20$ 秒。B：$29 - 16 = 13$ 秒。C：$34 - 19 =
15$ 秒。

\textbf{2.} 台站 B：其延迟最短，而延迟越短，
\omterm{def:g7:earthquakes:focus}{震中}越近。（其较快的波也最先到达。）

\textbf{3.} A：$8 \times 20 = \qty{160}{km}$。B：$8 \times 13 =
\qty{104}{km}$。C：$8 \times 15 = \qty{120}{km}$。

\textbf{4.} 到 A 的距离为 \qty{160}{km} 的所有点，构成半径为
\qty{160}{km}、以 A 为圆心的圆。\omterm{def:g7:earthquakes:focus}{震中}位于该圆上的某一点。

\textbf{5.} A：$160 \div 20 = 8$ 格。B：$104 \div 20 = 5.2$ 格。
C：$120 \div 20 = 6$ 格。

\textbf{6.} E 位于 O 以东 $3 + 6 = 9$ 格、以北 14.5 格处，
即三个台站之间的海上；A、B 两圆均经过该点。

\textbf{7.} \omterm{def:g7:earthquakes:focus}{震中}位于海底。海底的突然移动可能引发
\omterm{def:g7:earthquakes:tsunami}{海啸}：\omterm{def:g1:landscapes-around-us:water}{海岸}与岛屿都必须得到预警。

\textbf{8.} \omterm{def:g7:earthquakes:magnitude}{震级}每增加一级，能量约增大 32 倍。

\textbf{9.} $12 \times 20 = \qty{240}{km}$。

\textbf{10.} $240 \div 800 = 0.3$ 小时，即 $0.3 \times 60 = 18$ 分钟。

\textbf{11.} \omterm{def:g4:volcanoes-earthquakes-first-look:earthquake}{地震}发生后 5 分钟。

\textbf{12.} $18 - 5 = 13$：波特黑尔的居民约有
\textbf{13 分钟}的时间可以到达高地。（\omterm{def:g7:earthquakes:tsunami}{海啸}波在
岛屿附近的浅水中减速，这多出一点时间，但没有人应指望它：
他们应立即出发。）
\end{solution}
